\documentclass[10pt,a4paper]{amsart}
\usepackage{amsmath,amssymb,amsthm,graphicx}
% For commutative diagrams we use Xy-pic
\usepackage[all]{xy}

% We use 2cell for 2-commutative diagrams.
\xyoption{2cell}
% Theorem environments.
%
\theoremstyle{plain}
\newtheorem{theorem}[subsection]{Theorem}
\newtheorem{proposition}[subsection]{Proposition}
\newtheorem{lemma}[subsection]{Lemma}

\theoremstyle{definition}
\newtheorem{definition}[subsection]{Definition}
\newtheorem{example}[subsection]{Example}
\newtheorem{exercise}[subsection]{Exercise}
\newtheorem{situation}[subsection]{Situation}

\theoremstyle{remark}
\newtheorem{remark}[subsection]{Remark}
\newtheorem{remarks}[subsection]{Remarks}

\numberwithin{equation}{subsection}

% Macros
%
\def\lim{\mathop{\mathrm{lim}}\nolimits}
\def\colim{\mathop{\mathrm{colim}}\nolimits}
\def\Spec{\mathop{\mathrm{Spec}}}
\def\Hom{\mathop{\mathrm{Hom}}\nolimits}
\def\Ext{\mathop{\mathrm{Ext}}\nolimits}
\def\SheafHom{\mathop{\mathcal{H}\!\mathit{om}}\nolimits}
\def\SheafExt{\mathop{\mathcal{E}\!\mathit{xt}}\nolimits}
\def\Sch{\mathit{Sch}}
\def\Mor{\mathop{\mathrm{Mor}}\nolimits}
\def\Ob{\mathop{\mathrm{Ob}}\nolimits}
\def\Sh{\mathop{\mathit{Sh}}\nolimits}
\def\NL{\mathop{N\!L}\nolimits}
\def\CH{\mathop{\mathrm{CH}}\nolimits}
\def\proetale{{pro\text{-}\acute{e}tale}}
\def\etale{{\acute{e}tale}}
\def\QCoh{\mathit{QCoh}}
\def\Ker{\mathop{\mathrm{Ker}}}
\def\Im{\mathop{\mathrm{Im}}}
\def\Coker{\mathop{\mathrm{Coker}}}
\def\Coim{\mathop{\mathrm{Coim}}}

% Boxtimes
%
\DeclareMathSymbol{\boxtimes}{\mathbin}{AMSa}{"02}

%
% Macros for moduli stacks/spaces
%
\def\QCohstack{\mathcal{QC}\!\mathit{oh}}
\def\Cohstack{\mathcal{C}\!\mathit{oh}}
\def\Spacesstack{\mathcal{S}\!\mathit{paces}}
\def\Quotfunctor{\mathrm{Quot}}
\def\Hilbfunctor{\mathrm{Hilb}}
\def\Curvesstack{\mathcal{C}\!\mathit{urves}}
\def\Polarizedstack{\mathcal{P}\!\mathit{olarized}}
\def\Complexesstack{\mathcal{C}\!\mathit{omplexes}}
% \Pic is the operator that assigns to X its picard group, usage \Pic(X)
% \Picardstack_{X/B} denotes the Picard stack of X over B
% \Picardfunctor_{X/B} denotes the Picard functor of X over B
\def\Pic{\mathop{\mathrm{Pic}}\nolimits}
\def\Picardstack{\mathcal{P}\!\mathit{ic}}
\def\Picardfunctor{\mathrm{Pic}}
\def\Deformationcategory{\mathcal{D}\!\mathit{ef}}
\usepackage{cleveref}
\crefname{originallemma}{original lemma}{original lemmas}
\Crefname{originallemma}{Original lemma}{Original lemmas}
\makeatletter
\newlabel{divisors-lemma-blowing-up-gives-effective-Cartier-divisor}{{02OS}{source}{}{sourcepointer-1}{}}
\newlabel{divisors-lemma-blowing-up-gives-effective-Cartier-divisor@cref}{{[originallemma][0][]02OS}{[1][0][]source}{}{}{}}
\makeatother
\crefname{originallemma}{original lemma}{original lemmas}
\Crefname{originallemma}{Original lemma}{Original lemmas}
\makeatletter
\newlabel{divisors-lemma-blowing-up-projective}{{02NS}{source}{}{sourcepointer-2}{}}
\newlabel{divisors-lemma-blowing-up-projective@cref}{{[originallemma][0][]02NS}{[1][0][]source}{}{}{}}
\makeatother
\crefname{originallemma}{original lemma}{original lemmas}
\Crefname{originallemma}{Original lemma}{Original lemmas}
\makeatletter
\newlabel{divisors-lemma-blow-up-and-irreducible-components}{{0BFM}{source}{}{sourcepointer-3}{}}
\newlabel{divisors-lemma-blow-up-and-irreducible-components@cref}{{[originallemma][0][]0BFM}{[1][0][]source}{}{}{}}
\makeatother
\crefname{originallemma}{original lemma}{original lemmas}
\Crefname{originallemma}{Original lemma}{Original lemmas}
\makeatletter
\newlabel{varieties-lemma-finite-in-codim-1}{{0AB7}{source}{}{sourcepointer-4}{}}
\newlabel{varieties-lemma-finite-in-codim-1@cref}{{[originallemma][0][]0AB7}{[1][0][]source}{}{}{}}
\makeatother
\crefname{originallemma}{original lemma}{original lemmas}
\Crefname{originallemma}{Original lemma}{Original lemmas}
\makeatletter
\newlabel{algebra-lemma-locally-free-semi-local-free}{{02M9}{source}{}{sourcepointer-5}{}}
\newlabel{algebra-lemma-locally-free-semi-local-free@cref}{{[originallemma][0][]02M9}{[1][0][]source}{}{}{}}
\makeatother
\title{Two nonzerodivisors in a one-dimensional local ring acquire a common factor after a finite blowup}
\author{The Stacks Project Authors}
\date{stacks-project:chow.tex:local-snapshot-20261012; GNU FDL 1.2 or later}
\begin{document}
\maketitle
\noindent\textit{Editorial note.} All delivered mathematical text is the author's own, copied from the pinned Stacks Project chapter file chow.tex; no statement or proof step was rewritten, reordered or removed.\par\medskip
\noindent\textit{Editorial note.} source\_id and source\_version carry the current authoritative snapshot label of this source, 20261012, read from sources/stacks-chow/license-observation.json and SOURCE.json at the moment of opening.\par\medskip
\noindent\textit{Editorial note.} The preamble is the chapter's own native preamble excerpt used for this source: lines 14-18 (Xy-pic setup) and 154-217 (theorem declarations and mathematical macros) of sources/stacks-chow/preamble.tex. The quoted class line and the xr-hyper externaldocument block are not included because the assembler supplies a plain amsart wrapper; no macro was invented or redefined.\par\medskip
\noindent\textit{Editorial note.} All five referenced results live in earlier chapters, so their locators point at the verbatim lines of the project's own tag list staged in the snapshot (sources/stacks-chow/tags.txt) that carry the canonical tags.\par\medskip
\noindent\textit{Editorial note.} This is the first card using the new lane ident range L800-L999 (Chief 2026-10-12), which ends the cross-lane numbering collisions.\par\medskip
\noindent\textit{Allowed standard prerequisites (editorial).}\begin{itemize}
\item Blowing up in an ideal, and the effective Cartier divisor it produces. Divisors, Lemma labelled divisors-lemma-blowing-up-gives-effective-Cartier-divisor (canonical Stacks tag 02OS); invoked at proof lines 581-582 and again at 596. Supplies that the blowup X = Proj(oplus I\textasciicircum{}d) -> Spec(R) is an isomorphism over the punctured spectrum and that IR' is locally principal.
\item The blowup is projective. Divisors, Lemma labelled divisors-lemma-blowing-up-projective (canonical Stacks tag 02NS); invoked at proof lines 586-587. Supplies projectivity, which is combined with the next two inputs to conclude finiteness.
\item Generic points of the blowup and finiteness of a morphism finite in codimension one. Divisors, Lemma divisors-lemma-blow-up-and-irreducible-components (0BFM, line 589) and Varieties, Lemma varieties-lemma-finite-in-codim-1 (0AB7, line 590). Supply that every generic point of X lies over the punctured spectrum, so that the projective morphism is finite.
\item A locally free module of rank one over a semi-local ring is free. Algebra, Lemma labelled algebra-lemma-locally-free-semi-local-free (canonical Stacks tag 02M9); invoked at proof lines 598-599. Supplies that the locally principal ideal IR' is in fact principal, which produces the common factor t.
\end{itemize}\medskip
\ifdefined\hypertarget\hypertarget{sourcepointer-1}{}\fi
\noindent\textit{Source reference.} Original lemma 02OS: sources/stacks-chow/tags.txt, source lines 3527--3527. Cross-chapter standard prerequisite invoked twice, at proof lines 582 and 596, for the behaviour of the blowup over the punctured spectrum and local principality of IR'.\par\medskip
\ifdefined\hypertarget\hypertarget{sourcepointer-2}{}\fi
\noindent\textit{Source reference.} Original lemma 02NS: sources/stacks-chow/tags.txt, source lines 3491--3491. Cross-chapter standard prerequisite invoked once at proof line 587 for projectivity of the blowup.\par\medskip
\ifdefined\hypertarget\hypertarget{sourcepointer-3}{}\fi
\noindent\textit{Source reference.} Original lemma 0BFM: sources/stacks-chow/tags.txt, source lines 14390--14390. Cross-chapter standard prerequisite invoked once at proof line 589 for the location of the generic points of the blowup.\par\medskip
\ifdefined\hypertarget\hypertarget{sourcepointer-4}{}\fi
\noindent\textit{Source reference.} Original lemma 0AB7: sources/stacks-chow/tags.txt, source lines 13010--13010. Cross-chapter standard prerequisite invoked once at proof line 590 for the finiteness criterion that makes the blowup affine.\par\medskip
\ifdefined\hypertarget\hypertarget{sourcepointer-5}{}\fi
\noindent\textit{Source reference.} Original lemma 02M9: sources/stacks-chow/tags.txt, source lines 3436--3436. Cross-chapter standard prerequisite invoked once at proof lines 598-599 to upgrade locally principal to principal over the semi-local ring R'; the final writing a = a't, b = b't is retained in the counted span.\par\medskip
\begin{lemma}
\label{lemma-Noetherian-domain-dim-1-two-elements}
Let $(R, \mathfrak m)$ be a Noetherian local ring of dimension $1$.
Let $a, b \in R$ be nonzerodivisors.
There exists a finite ring extension $R \subset R'$
with $R'/R$ annihilated by a power of $\mathfrak m$
and nonzerodivisors $t, a', b' \in R'$ such that
$a = ta'$ and $b = tb'$ and $R' = a'R' + b'R'$.
\end{lemma}

\begin{proof}
If $a$ or $b$ is a unit, then the lemma is true with $R = R'$.
Thus we may assume $a, b \in \mathfrak m$.
Set $I = (a, b)$. The idea is to blow up $R$ in $I$.
Instead of doing the algebraic argument we work geometrically.
Let $X = \text{Proj}(\bigoplus_{d \geq 0} I^d)$.
By Divisors, Lemma
\ref{divisors-lemma-blowing-up-gives-effective-Cartier-divisor}
the morphism $X \to \Spec(R)$ is an isomorphism over
the punctured spectrum $U = \Spec(R) \setminus \{\mathfrak m\}$.
Thus we may and do view $U$ as an open subscheme of $X$.
The morphism $X \to \Spec(R)$ is projective by
Divisors, Lemma \ref{divisors-lemma-blowing-up-projective}.
Also, every generic point of $X$ lies in $U$, for example
by Divisors, Lemma \ref{divisors-lemma-blow-up-and-irreducible-components}.
It follows from Varieties, Lemma \ref{varieties-lemma-finite-in-codim-1}
that $X \to \Spec(R)$ is finite. Thus $X = \Spec(R')$ is
affine and $R \to R'$ is finite. We have $R_a \cong R'_a$ as $U = D(a)$.
Hence a power of $a$ annihilates the finite $R$-module $R'/R$.
As $\mathfrak m = \sqrt{(a)}$ we see that $R'/R$ is annihilated
by a power of $\mathfrak m$. By
Divisors, Lemma \ref{divisors-lemma-blowing-up-gives-effective-Cartier-divisor}
we see that $IR'$ is a locally principal ideal.
Since $R'$ is semi-local we see that $IR'$ is principal,
see Algebra, Lemma \ref{algebra-lemma-locally-free-semi-local-free},
say $IR' = (t)$. Then we have $a = a't$ and $b = b't$ and everything is
clear.
\end{proof}

\end{document}
